\section{Compressed least representatives from an orbit trace}
\label{sec:transversal}

An orbit count alone cannot reveal which boundary circles lie in a particular
surface component. We therefore compile a stronger output from the existing
trace: one original point of each orbit, still encoded as intervals.

\begin{definition}
For an equivalence relation $\sim$ on $[0,N)\cap\Z$, its least transversal is
\[
 R(\sim)=\{\min O:O\text{ is a }\sim\text{-class}\}.
\]
An interval representation of $R$ is its canonical union of nonempty,
pairwise disjoint and nonadjacent half-open integer intervals.
\end{definition}

An output interval $[a,b)$ means that each of its $b-a$ integer points is a
representative of a different orbit. It does not mean that those points form
one orbit. For a relation with no pairings, the complete answer is the single
interval $[0,N)$, even when $N$ is enormous.

\begin{theorem}[Least-transversal compiler]
\label{thm:transversal}
Let a complete valid maintained AHT trace have $E$ local events and a total
of $G$ static gap intervals. Its least original transversal can be computed
as at most $G$ half-open intervals. The producer and an independently
implemented verifier take time polynomial in $E$, $G$, and
$\log(N+1)$, in addition to the inherited cost of validating the trace.
Their output is independent of which valid complete AHT schedule was used.
For $N=0$, the output is empty.
\end{theorem}

\subsection{Forward construction and invariant}

Maintain an increasing injection $\iota$ from the current ordered point
universe into the original universe. Store its image as a sorted list of
original intervals. Concatenating the lengths of these intervals gives
current ranks, so a current interval can be pulled back to original labels
by splitting only at the interval-list boundaries.

Initially the live list is $[0,N)$ and nothing has been emitted. We maintain
three assertions: emitted points are the least points of distinct completed
original orbits; every uncompleted original orbit has live points, including
its least point; and the current relation, transported by $\iota$, is the
restriction of the original relation to these live points.

Identity deletion, trimming, merging, and transmission preserve the point
universe and its orbit relation. They leave the live list unchanged.
Only contractions and truncations require additional work.

A certified suffix truncation retains the prefix $[0,c)$ of the current
universe. Every deleted point is equivalent to a retained point. For a
translation carrier, iterated inverse translation reaches that prefix; the
operation computes this by division, not iteration over the number of
periods. For a trimmed reflection, the relevant domain precedes the deleted
image, and one reflection reaches it. The inherited local preconditions
also ensure that the retained relation is the induced relation on the
prefix. As $\iota$ is increasing, every deleted point has a smaller original
point in its orbit that survives. No uncompleted orbit loses its minimum.
Replace the live list by the intervals representing its first $c$ points.

A static contraction removes uncovered current intervals. Each point there
is a singleton of the current relation. By the invariant it is the sole live
point of one original orbit, and that point is its minimum. Intersect the
current gaps with the live interval list, emit the corresponding original
interval pieces, delete those pieces, and close the current-rank gaps.
An original point is emitted at most once. When the trace ends, no live
points remain, so the emitted union is exactly $R(\sim)$. This proves the
semantic assertion and schedule independence in \cref{thm:transversal}.

\subsection{A different reverse construction}

The certificate checker first replays the source-bound AHT proof with the
existing independent checker. It then reconstructs the selected set
\emph{backwards}, starting with the empty final universe. Undoing a
truncation appends an unselected suffix. Undoing a contraction reinserts
every static gap as selected points and lifts the previously selected
points across those insertions. Other events leave point indices unchanged.
This method uses neither forward live-interval tracking nor orbit discovery.

Here is an explicit formula. Let the half-open gaps before a contraction be
$[a_i,b_i)$ in increasing order. Put
\[
 \ell_i=b_i-a_i,\qquad
 c_i=a_i-\sum_{j<i}\ell_j.
\]
The contracted position of gap $i$ is $c_i$. A surviving contracted point
$q$ lifts to
\begin{equation}
 \lambda(q)=q+\sum_{i:c_i\le q}\ell_i.
 \label{eq:gap-lift}
\end{equation}
For an existing selected interval $[u,v)$, add the hull
\[
 [\lambda(u),\lambda(v-1)+1)
\]
and add every reinserted gap $[a_i,b_i)$. The extra points introduced by the
hull are exclusively gap points, which are themselves selected. Thus the
union is exact. Prefix sums and binary searches evaluate \eqref{eq:gap-lift}
without visiting any represented point.

\begin{lemma}[Interval bound]
\label{lem:selector-size}
The final selected set has at most $G$ intervals. During forward compilation
there are at most $1+G$ live original intervals, and at most $1+2G$ raw
emitted interval pieces in total.
\end{lemma}
\begin{proof}
In the reverse construction, inserting $g$ selected gaps increases the
number of selected intervals by at most $g$. An existing selected interval,
together with the gaps crossed by its lift, still has a single selected
hull. Additional isolated gaps account for at most $g$ new intervals.
Reverse truncation adds no selected interval. Starting with zero proves
the bound $G$.

For the live-list bound, deleting one current gap can split at most one
live original interval into two retained pieces; deleting a suffix cannot
increase the count. For raw emissions, splitting a live list of $p$
intervals at the $2g$ endpoints of a contraction yields at most $p+2g$
pieces. Each piece is retained or emitted. Summing this inequality over
contractions telescopes the live-list counts; suffix deletions only decrease
them. The initial list has at most one interval, giving $1+2G$ raw emissions.
Sorting and coalescing do not increase their number.
\end{proof}

\subsection{Costs and certificate obligations}

A straightforward list implementation takes
\[
 O\bigl(E(1+G)+G\log(G+1)\bigr)
\]
endpoint operations for the forward construction. A conservative bound for
the reverse selector is
\[
 O\bigl((E+G)(G+1)\log(G+2)\bigr)
\]
endpoint operations, besides inherited orbit-proof replay. All endpoints
have $O(\log(N+1))$ bits. These bounds establish polynomial encoded cost;
they are not claimed to be degree-optimal. The displayed dependence on $G$
also identifies where a persistent interval tree could improve the current
list implementation.

The certificate contains the exact original size, its orbit proof, the
claimed orbit count, and canonical representative intervals. The verifier
checks the orbit proof against the caller's actual pairing list, reconstructs
the reverse selector, and compares the entire interval list and its
cardinality. Matching cardinality alone would be unsound: a different
transversal can have the correct number of points but fail the least-point
claim, and an arbitrary set of that size need not be a transversal at all.

The discovery cycle allowance applies to the inherited discovery phase.
Trace compilation and later independent checking use cooperative cancellation
and explicit trace size as their resource parameters. Incomplete discovery
produces neither a transversal certificate nor a purported completed count.

\section{Turning nested component counts into additive weights}
\label{sec:nested}

\begin{theorem}[Fine components inside coarse components]
\label{thm:nested}
Let $\sim_f$ and $\sim_c$ be finite interval-encoded equivalence relations.
Suppose an injective interval map $i$ from the fine universe into the coarse
universe sends each fine orbit into a single coarse orbit. If $R_f$ is a
transversal for $\sim_f$, then the weight $\mathbf1_{i(R_f)}$ has sum
\[
 \sum_{p\in O}\mathbf1_{i(R_f)}(p)
 =\#\{A\text{ a fine orbit}:i(A)\subseteq O\}
\]
on every coarse orbit $O$.
\end{theorem}
\begin{proof}
Every fine orbit contributes exactly its unique representative, and that
representative lies in the same coarse orbit as all its fine-orbit points.
Distinct fine representatives remain distinct under the injection. Thus
the contributions are exactly one for each contained fine orbit and zero
otherwise.
\end{proof}

The interval encoding of the injection is a hypothesis, not something that
can be inferred from two lists of component counts. In the application,
the fine relation is the boundary-arc relation, the coarse relation is the
full surface relation, and the injection is the literal inclusion of
boundary edge points. All these data are reconstructed from the original
validated triangulation and vector.

\subsection{Boundary markers in the full edge order}

Boundary points are grouped by boundary edge in increasing global edge
order. Full surface points are grouped by all edges in that same order.
For a boundary edge $e$, let $o_\partial(e)$ and $o_F(e)$ be its two block
offsets, and let $n_e$ be its intersection count. The inclusion on its block is
\[
 [o_\partial(e),o_\partial(e)+n_e)\longrightarrow
 [o_F(e),o_F(e)+n_e),\qquad
 p\longmapsto p+o_F(e)-o_\partial(e).
\]
Both orders use the same global edge orientation, so no further reflection
is needed. A selector interval is split where it crosses boundary-edge
block endpoints and each piece is translated by this formula.

If the boundary selector has $r$ intervals and there are $e_\partial$
boundary-edge blocks, the total number of pieces is $O(r+e_\partial)$.
This follows by merging two sorted endpoint lists; it is not a product
bound $r e_\partial$. Since $r\le G_\partial$ and
$e_\partial=O(t)$, the embedded markers remain polynomially encoded.

\subsection{Integral Euler ownership}

Give each surface vertex weight $+1$, place weight $-1$ for each surface
arc at one of its endpoints, and place weight $+1$ for each normal disc at
one chosen corner. Every charge is located within its cell's surface
component. Thus the orbit sum is the Euler characteristic of that component.

This ownership is compressed explicitly. The full edge-point universe
receives $+1$. The source interval of each original face-arc pairing
receives $-1$, with its inclusive endpoint changed to a half-open stop.
For each nonempty normal disc type, the chosen corners are consecutive
along a suitable tetrahedron edge, so one interval carries their $+1$
charges. The original arc list is used for these cell charges, before
deleting identity pairings or otherwise reducing the orbit relation.
Deleting a generator is not a license to omit a geometric edge charge.

Combine the Euler weight with the boundary-marker weight:
\begin{equation}
 w(p)=\bigl(w_\chi(p),\mathbf1_{i(R_\partial)}(p)\bigr)\in\Z^2.
 \label{eq:two-weights}
\end{equation}
By \cref{thm:nested}, the weighted orbit sum on a connected component $C$ is
exactly $(\chi(C),b(C))$. In particular, a closed component has second
coordinate zero, and an annulus has second coordinate two regardless of
how many normal arcs subdivide either boundary circle.
